\section{The full ordinary Dirichlet-to-Neumann map}
\label{sec:full-conclusion}
\begin{proof}[Proof of Theorem~\ref{main:full}]
The common tangential metric gives $dS_{g_1}=dS_{g_2}=dS$.
For smooth boundary functions Green's formula gives
\[
 Q_{g_j}(f,h)=\int_{\partial M}(\Lambda_{g_j}f)\bar h\,dS.
\]
Equality of the ordinary maps therefore gives equality of the energy
forms, first on smooth functions and then on $H^{1/2}$ by their boundedness.
Choose a boundary point $p$ and set $\Gamma=\partial M\setminus\{p\}$.
Since $\dim\partial M=n-1\ge2$, this is nonempty, proper, relatively
open, and dense, including when the boundary is disconnected.
Theorem~\ref{main:patch} supplies a smooth $\Phi$ with
$g_2=\Phi^*g_1$ fixing $\Gamma$. Continuity fixes the remaining point.
Taking $F=\Phi^{-1}$ gives $g_1=F^*g_2$ and fixes every boundary point,
as required.
\end{proof}
The data used are exactly one full, ordinary, zero-frequency map.
All dimensions $n\ge3$, arbitrary smooth positive metrics, nonorientable
manifolds, and arbitrary boundary components have been retained.
